%==============================================================
% Advanced Scientific Writing -- Assignment 3
% Position-paper OUTLINE, typeset with the official PSB / World
% Scientific author kit (ws-procs11x85 document class).
%
% Build:  pdflatex twice  (or: latexmk -pdf -cd)
%
% OUTLINE: each section gives its heading plus short bullets
% describing the content that section will contain.
%==============================================================

\documentclass{ws-procs11x85}
\usepackage{ws-procs-thm}   % comment out if amsthm/theorem/ntheorem is used

\begin{document}

\title{Nicotinic Acetylcholine Receptor Prediction using VEP for Humans}

\author{Hari Krishna Mohan Raj}

\address{Department of Computational Chemistry,\\
Hochschule Bonn Rhein Sieg, Sankt Augustin, Germany\\
E-mail: harikrishnam0711@gmail.com}

\begin{abstract}
Variants in human nicotinic acetylcholine receptor (nAChR) genes cause
a range of neurological and neuromuscular disorders, yet most newly
observed variants are of unknown consequence. Existing predictors
report only whether a variant is damaging, not whether it causes
loss of function (LOF) or gain of function (GOF) a distinction that
is decisive for therapy. We argue that variant effect prediction for
nAChRs must move beyond binary pathogenicity toward receptor-specific,
direction-aware prediction, and we outline the case, the gap in
current tools, and what such a predictor would require.
\end{abstract}

\keywords{Variant effect prediction; nicotinic acetylcholine receptor;
gain of function; loss of function; ion channel; machine learning.}

% required, do-not-remove
\copyrightinfo{\copyright\ 2026 Hari Krishna Mohan Raj, Karl Kirschner. Open Access chapter
published by World Scientific Publishing Company and distributed under
the terms of the Creative Commons Attribution Non-Commercial (CC
BY-NC) 4.0 License.}

%==============================================================
\section{Introduction}\label{sec:intro}
%==============================================================
\begin{itemlist}
\item nAChR variants cause myasthenic syndromes and epilepsy.
\item sequencing outpaces functional assays; most variants are VUS.
\item Core point: ``pathogenic'' is not equal to ``directional'' LOF and GOF need opposite therapies.
\item nAChR prediction must become direction-aware and receptor-specific.
% \item Roadmap of the paper.
\end{itemlist}

%==============================================================
\section{Background: nAChRs and the Weight of Direction}\label{sec:background}
%==============================================================
\begin{itemlist}
\item Brief nAChR primer, pentameric ligand gated channels, muscle vs neuronal types.
\item Genotype to phenotype, which subunit genes map to which disorders.
% \item Why direction, not just pathogenicity, drives treatment.
\item LOF/GOF is measurable experimentally, but slow and costly.
\end{itemlist}

%==============================================================
\section{The Current Landscape of Variant Effect Prediction}\label{sec:landscape}
%==============================================================
\begin{itemlist}
\item Talk about what current tools can do and cant. 
\end{itemlist}


%==============================================================
\section{The Gap: Why Existing Tools Fall Short for nAChRs}\label{sec:gap}
%==============================================================
\begin{itemlist}
\item No nAChR specialized predictor; subunit identity ignored.
\item Direction tools target other channels Na/Ca, thin nAChR coverage.
\item nAChR structural determinants (binding interface, pore, gating) unused.
\item Data scarce
\end{itemlist}


%==============================================================
\section{Challenges and Open Questions}\label{sec:challenges}
%==============================================================
\begin{itemlist}
\item Data scarcity and curation burden.
\item LOF/GOF as a spectrum, mixed effect variants.
\item Generalization across subunits and to novel variants.
\item Interpretability and the risk of overclaiming.
\item No shared nAChR direction of effect benchmark.
\end{itemlist}

%==============================================================
\section{Conclusion and Call to Action}\label{sec:conclusion}
%==============================================================
\begin{itemlist}
\item Conclude the results
\item talk about faster, cheaper triage of nAChR VUS.
\end{itemlist}


%==============================================================
% References
%==============================================================
\begin{thebibliography}{9}

\bibitem{sift} P.~C. Ng and S.~Henikoff, SIFT: predicting amino acid
changes that affect protein function, {\em Nucleic Acids Res.}
{\bf 31}, 3812 (2003).

\bibitem{polyphen2} I.~A. Adzhubei {\em et al.}, A method and server
for predicting damaging missense mutations, {\em Nat. Methods}
{\bf 7}, 248 (2010).

\bibitem{cadd} M.~Kircher {\em et al.}, A general framework for
estimating the relative pathogenicity of human genetic variants,
{\em Nat. Genet.} {\bf 46}, 310 (2014).

\bibitem{alphamissense} J.~Cheng {\em et al.}, Accurate proteome-wide
missense variant effect prediction with AlphaMissense, {\em Science}
{\bf 381}, eadg7492 (2023).

\bibitem{esm1b} N.~Brandes {\em et al.}, Genome-wide prediction of
disease variant effects with a deep protein language model,
{\em Nat. Genet.} {\bf 55}, 1512 (2023).

\bibitem{funncion} A.~Brunklaus {\em et al.}, Gene variant effects
across sodium channelopathies predict function and guide precision
therapy, {\em Brain} {\bf 145}, 4275 (2022).

\bibitem{logofunc} D.~Stein {\em et al.}, Genome-wide prediction of
pathogenic gain- and loss-of-function variants from ensemble learning
of a diverse feature set, {\em Genome Med.} {\bf 15}, 103 (2023).

\bibitem{premode} G.~Zhang {\em et al.}, Predicting the mode of action
of missense variants (PreMode), bioRxiv (2024).

\bibitem{mohanraj2026} H.~K. Mohan Raj, {\em Variant Effect Predictor
in Human nAChR Proteins: Abstract and Publishable Results}, Advanced
Scientific Writing, Assignment~2 (2026).


\end{thebibliography}

\end{document}
